% Generated draft: requires mathematical and source review.
\documentclass[11pt]{article}
\usepackage[margin=25mm]{geometry}
\usepackage{fontspec}
\setmainfont{Latin Modern Roman}
\usepackage{amsmath,amssymb,mathtools,unicode-math}
\setmathfont{Latin Modern Math}
\usepackage{enumitem,needspace,xurl,hyperref}
\setlength{\parindent}{0pt}
\setlength{\parskip}{5pt}
\setlength{\emergencystretch}{3em}
\newcommand{\N}{\mathbb{N}}
\newcommand{\Z}{\mathbb{Z}}
\newcommand{\Q}{\mathbb{Q}}
\newcommand{\R}{\mathbb{R}}
\newcommand{\C}{\mathbb{C}}
\newcommand{\F}{\mathbb{F}}
\newcommand{\A}{\mathbb{A}}
\newcommand{\PP}{\mathbb{P}}
\newcommand{\E}{\mathbb{E}}
\newcommand{\eps}{\varepsilon}
\newcommand{\dd}{\,\mathrm{d}}
\begin{document}
\section*{UnsolvedMath ID: 20000389}
\textbf{Source problem code:} catalogue entry
\subsection{Definitions and mathematical statement}
John Voight asks when counting formulas become simpler after allowing a partial compactification. A partial compactification is an open immersion $U\hookrightarrow Y$ adding specified boundary strata; $Y$ need not itself be proper. A precise counting question must choose the spaces, allowed boundary, finite fields, ordinary versus stack-weighted counts, and a meaning of ``nicer.'' None is fixed in the short original.

The attached unreviewed aid gives two models. For positive integer weights, let $\PP(w_0,\ldots,w_m)=[(\A^{m+1}\setminus0)/\mathbb G_m]$. Its groupoid mass over $\F_q$ is the sum of $1/|\operatorname{Aut}(x)(\F_q)|$ over isomorphism classes of rational objects. The aid reports
\[
\#\PP(w_0,\ldots,w_m)(\F_q)=\frac{q^{m+1}-1}{q-1}.
\]
It consequently reports, in characteristic greater than three, mass $q+1$ for the compactified elliptic stack, mass $1$ for its nodal boundary and mass $q$ for the smooth open. These are weighted stack counts, not counts of isomorphism classes with weight one.

For a smooth finite-type curve $C/\F_q$, define $Z_C(t)=\prod_{x\in|C|}(1-t^{\deg x})^{-1}$. Let $a_{n,r}$ count effective degree-$n$ divisors whose multiplicities at every closed point are at most $r$, where $r\ge1$. The aid reports
\[
\sum_{n\ge0}a_{n,r}t^n=\frac{Z_C(t)}{Z_C(t^{r+1})}.
\]
For $C=\A^1$ this gives $a_{n,r}=q^n-q^{n-r}$ for $n\ge r+1$. Thus in monic polynomial coefficient space, for $n\ge3$, replacing the squarefree condition by allowing double roots but no triple roots changes the count from $q^n-q^{n-1}$ to $q^n-q^{n-2}$. The added stratum has $q^{n-1}-q^{n-2}$ coefficient points, and the omitted triple-root locus has codimension two. For equations $y^2=f(x)$, a nodal interpretation additionally needs characteristic not two and suitable curve/compactification conventions; these coefficient counts are not moduli masses or a count of individual nodes.

The aid suggests that independent closed-point factors or simple quotient-stack masses can make boundary corrections explicit. This is a sufficient mechanism illustrated by the two models, not a general characterization of all beneficial partial compactifications.
\subsection{Short English statement}
Identify partial compactifications whose added boundary gives a simpler arithmetic count. The aid supplies weighted elliptic-stack and bounded-root-multiplicity examples, while the general notion of a simpler count remains unspecified.
\subsection{Sources}
\begin{itemize}
\item John Voight, Question 11, AIM \emph{Arithmetic statistics over finite fields and function fields}: \url{https://aimath.org/pastworkshops/arithstatffieldproblems.pdf}. Canonical location: \texttt{aim-algebraic-number-theory-notes.json}, notes[40]. The original is labeled exact; remarks and literature are empty. The supplied extraction is used without consulting a PDF.
\item The attached boundary-correction aid is machine generated, not peer reviewed, with low novelty confidence. It supplies no complete authoritative references for the stack or divisor formulas.
\item Local registry: \texttt{data/unsolvedmath/problems.json}, permanent catalogue entry, exact source code \texttt{catalogue entry}.
\end{itemize}
\end{document}
